\documentclass[conference]{IEEEtran}

\usepackage{graphicx}
\usepackage{amsmath,amssymb}
\usepackage{float}
\usepackage{booktabs}
\usepackage{multirow}

\begin{document}

\title{
    \textbf{Benchmark Report: 526.blender\_r} \\
    \text{Register Spill Analysis on RISC-V via gem5}
}

\author{
    \IEEEauthorblockN{Lütfullah Burak Kaya}
    \IEEEauthorblockA{
        Ozyegin University \\
        burak.kaya.50711@ozu.edu.tr
    }
    \and
    \IEEEauthorblockN{Ismail Akturk}
    \IEEEauthorblockA{
        Ozyegin University \\
        ismail.akturk@ozyegin.edu.tr
    }
}

\newcommand{\LongDate}{%
    \number\day\space
    \ifcase\month
    \or January\or February\or March\or April\or May\or June%
    \or July\or August\or September\or October\or November\or December%
    \fi
    \space \number\year
}

\twocolumn[
\begin{@twocolumnfalse}
\maketitle
\begin{center}{\small \LongDate}\end{center}
\vspace{1em}
\end{@twocolumnfalse}
]

% ================================================

\begin{abstract}
This report presents the register spill characterization of the
\texttt{526.blender\_r} benchmark from SPEC CPU2017, executed on a RISC-V
(RV64GC) target using the gem5 TimingSimpleCPU simulator.
A custom SpillDetector was integrated into gem5 to count store-then-load
patterns on the stack within a user-defined Region of Interest (ROI).
The test dataset (\texttt{cube.blend}, 1 frame) was used as the input.
Results show a spill rate of approximately 6.22\% within the ROI ---
the highest observed across all SPEC CPU2017 benchmarks evaluated in
this study --- corresponding to over 61 million detected spill events
out of 994 million ROI instructions.
ROI instrumentation required a non-trivial fix: the original markers
were placed in a dead code path (\texttt{cycles\_standalone.cpp}) and
had to be moved to the actual render entry point in \texttt{creator.c}.
\end{abstract}

% =========================================================
\section{Benchmark Overview}

\texttt{526.blender\_r} is the SPEC CPU2017 rate variant of Blender,
the open-source 3D rendering suite. The benchmark exercises Blender's
ray-tracing renderer by rendering one or more frames of a 3D scene.
It is compute-intensive and mixed-precision, combining floating-point
geometry calculations with integer-heavy BVH tree traversal and
pixel compositing.

The RISC-V binary (\texttt{blender\_r.riscv}) was compiled with gem5 ROI
markers inserted in \texttt{creator.c}, around the \texttt{RE\_BlenderAnim()}
call inside the \texttt{render\_animation()} callback:

\begin{itemize}
    \item \texttt{m5\_work\_begin(0,0)} --- before \texttt{RE\_BlenderAnim()}
    \item \texttt{m5\_work\_end(0,0)}   --- after  \texttt{RE\_BlenderAnim()}
\end{itemize}

Only spills occurring between these two markers are counted.

\subsection{ROI Instrumentation History}

An earlier version of the binary placed ROI markers inside
\texttt{cycles\_standalone.cpp}'s \texttt{main()} function. That
function is the entry point for a \emph{standalone} Cycles renderer
binary and is \emph{never called} when running \texttt{blender\_r},
whose entry point resides in \texttt{creator.c}. As a result, the
early binary produced no \texttt{riscv\_spill\_stats.txt} output.
The binary was rebuilt after moving the markers to the correct location.

% =========================================================
\section{Simulation Setup}

\subsection{Input Datasets}

Blender's three standard SPEC input sizes are summarised in
Table~\ref{tab:inputs}.

\begin{table}[H]
\centering
\caption{Blender Input Dataset Sizes}
\label{tab:inputs}
\begin{tabular}{llcc}
\toprule
\textbf{Dataset} & \textbf{Scene file} & \textbf{Frames} & \textbf{File size} \\
\midrule
test    & cube.blend          & 1 (frame 1)   & 437\,KB \\
train   & sh5\_reduced.blend  & 1 (frame 234) & 41\,MB  \\
refrate & sh3\_no\_char.blend & 1 (frame 849) & 57\,MB  \\
\bottomrule
\end{tabular}
\end{table}

The \textbf{test} dataset (\texttt{cube.blend}) was selected for this run
to obtain a fast and verifiable measurement consistent with the policy
used for other benchmarks in this study.

\subsection{gem5 Configuration}

\begin{itemize}
    \item \textbf{CPU model:}   TimingSimpleCPU
    \item \textbf{ISA:}         RISC-V RV64GC
    \item \textbf{Caches:}      enabled (default L1 I/D + L2)
    \item \textbf{Memory:}      4\,GB simulated physical RAM
    \item \textbf{Spill mode:}  summary only (\texttt{spill\_verbose=False})
\end{itemize}

The full gem5 invocation was:

\begin{verbatim}
./build/RISCV/gem5.opt
  --outdir=benchmarks/blender/m5out_test2
  configs/deprecated/example/se.py
  --cmd=benchmarks/blender/blender_r.riscv
  --options="benchmarks/blender/cube.blend
             --render-output benchmarks/blender/cube_
             --threads 1 -b -F RAWTGA -s 1 -e 1 -a"
  --cpu-type=TimingSimpleCPU
  --caches
  --mem-size=4GB
\end{verbatim}

% =========================================================
\section{Simulation Results}

The simulation completed successfully. Blender printed the expected
frame-by-frame progress log and exited cleanly
(\texttt{Blender quit}).

\subsection{Pre-ROI Context}

Before the ROI was entered, the binary executed Blender's startup,
scene parsing, and renderer initialisation. Table~\ref{tab:preroi}
shows the global counters at the moment \texttt{m5\_work\_begin}
was triggered.

\begin{table}[H]
\centering
\caption{Global Counters at ROI Entry}
\label{tab:preroi}
\begin{tabular}{lr}
\toprule
\textbf{Metric} & \textbf{Value} \\
\midrule
Total instructions (pre-ROI) & 189,437,743 \\
Total spills detected (pre-ROI) & 6,449,141 \\
\bottomrule
\end{tabular}
\end{table}

\subsection{ROI Spill Statistics}

Table~\ref{tab:roi} presents the spill statistics collected exclusively
within the ROI (rendering of frame 1 of \texttt{cube.blend}).

\begin{table}[H]
\centering
\caption{ROI Spill Statistics --- 526.blender\_r (test, 1 frame)}
\label{tab:roi}
\begin{tabular}{lr}
\toprule
\textbf{Metric} & \textbf{Value} \\
\midrule
ROI instructions         &   994,414,074 \\
ROI stores               &   116,509,594 \\
ROI loads                &   233,612,215 \\
\midrule
\textbf{ROI spills}      & \textbf{61,845,869} \\
\textbf{Spill rate}      & \textbf{6.2193\%} \\
\bottomrule
\end{tabular}
\end{table}

\subsection{Timing}

\begin{table}[H]
\centering
\caption{Simulation Timing}
\label{tab:timing}
\begin{tabular}{lr}
\toprule
\textbf{Metric} & \textbf{Value} \\
\midrule
Simulated time  & 2.278\,s \\
Host wall time  & 842.72\,s ($\approx$14\,min) \\
Total simulated instructions & 1,191,956,028 \\
\bottomrule
\end{tabular}
\end{table}

% =========================================================
\section{Analysis}

\subsection{Spill Rate Interpretation}

The measured spill rate of \textbf{6.22\%} means that approximately
1 in every 16 instructions inside the ROI is a register spill
(a stack load matching a prior stack store). This is the highest
spill rate observed among all SPEC CPU2017 benchmarks evaluated
in this study, as shown in Table~\ref{tab:compare}.

\begin{table}[H]
\centering
\caption{Cross-Benchmark Spill Rate Comparison (test input)}
\label{tab:compare}
\begin{tabular}{lrr}
\toprule
\textbf{Benchmark} & \textbf{ROI Spills} & \textbf{Spill Rate} \\
\midrule
526.blender\_r & 61,845,869  & \textbf{6.2193\%} \\
544.nab\_r     & 151,310,616 & 2.2612\% \\
505.mcf\_r     & 446,113,136 & 1.8263\% \\
508.namd\_r    & 327,718,955 & 1.0269\% \\
\bottomrule
\end{tabular}
\end{table}

\subsection{Store vs.\ Load Balance}

The load count (233.6M) is approximately $2.0\times$ higher than the
store count (116.5M). This 2:1 ratio is notably lower than the 4.3:1
ratio seen in NAMD, suggesting that Blender's rendering kernel
both writes and reads the stack more symmetrically --- consistent
with the recursive BVH traversal pattern typical of ray tracers,
where activation frames are pushed and popped in nearly equal measure.

\subsection{Spill Fraction of Loads}

Within the ROI, the ratio of spills to total loads is:
\[
\frac{61{,}845{,}869}{233{,}612{,}215} \approx 26.5\%
\]
meaning roughly 1 in 4 load instructions is a spill reload.
This is an exceptionally high fraction compared to other
benchmarks and reflects the deep call-stack depth and high
register pressure of the recursive ray-tracing engine.

\subsection{Pre-ROI vs.\ ROI Spill Density}

Pre-ROI spills: $6{,}449{,}141$ over $189{,}437{,}743$ instructions
= \textbf{3.40\%} spill rate. This is notably lower than the ROI
rate (6.22\%), which is the inverse of the pattern seen in most
other benchmarks. In Blender, the core rendering phase is
\emph{more} register-pressure-intensive than startup, reflecting
the deep recursion and many live variables in the BVH traversal
and shading kernels during actual rendering.

\subsection{ROI Fraction of Total Execution}

The ROI accounts for $994{,}414{,}074$ out of $1{,}191{,}956{,}028$
total simulated instructions:
\[
\frac{994{,}414{,}074}{1{,}191{,}956{,}028} \approx 83.4\%
\]
This confirms that the ROI captures the dominant portion of the
benchmark's execution and is a faithful representative of its
overall behaviour.

% =========================================================
\section{Conclusion}

The \texttt{526.blender\_r} benchmark exhibits the highest register
spill rate (6.22\%) among all SPEC CPU2017 benchmarks evaluated in
this study. With 61.8 million spill events across 994 million ROI
instructions for a single rendered frame, Blender places severe
register pressure on the RISC-V backend. The 2:1 load-to-store ratio
and the 26.5\% spill fraction of loads both confirm that ray-tracing
recursion and BVH traversal generate dense stack traffic. These
results suggest that \texttt{526.blender\_r} would be a prime
candidate for investigating compiler spill-reduction strategies or
hardware register-file extensions on RISC-V.

\end{document}
